\chapter{Filtro pasa banda}
\label{chap:pasa-banda}

\section{Especificaciones}
\label{sec:pbd-especificaciones}

\begin{table}[!ht]
  \centering
  \begin{tabular}{|l|c|}
    \hline
    Parámetro & Valor \\
    \hline
    Orden & 2 \\
    Atenuación máxima en banda de paso $A_{max}$ & $\SI{0.1}{\decibel}$ \\
    Frecuencia central $f_o$ & $\SI{28}{\kilo\hertz}$ \\
    Ancho de banda $BW = f_{p2} - f_{p1}$ & $\SI{1.7}{\kilo\hertz}$ \\
    Ganancia en banda de paso & 9 V/V \\
    Aproximación & Butterworth \\
    Topología & Sallen-Key \\
    \hline
  \end{tabular}
  \caption{Especificaciones del filtro pasa banda.}
  \label{tab:pbd-especificaciones}
\end{table}

\section{Diseño}
\label{sec:pbd-diseno}

\subsection{Factor de calidad y función de transferencia}
\label{subsec:pbd-q}

Un pasa banda de segundo orden queda definido por su frecuencia central y su factor de calidad

\begin{equation}
  Q = \frac{f_o}{BW}
  \label{eq:pbd-q}
\end{equation}

y su función de transferencia es

\begin{equation}
  H(s) = \frac{A_m \, \dfrac{\omega_0}{Q} \, s}{s^2 + \dfrac{\omega_0}{Q} s + \omega_0^2}
  \label{eq:pbd-hs}
\end{equation}

% TODO: calcular Q y discutir la sensibilidad de la topología Sallen-Key para Q elevados (la ganancia interna K se
% acerca a 3 y el circuito se vuelve muy sensible a las tolerancias). Relacionar A_max con el ancho de banda definido
% a -A_max dB en lugar de -3 dB si corresponde.

\subsection{Cálculo de la etapa Sallen-Key}
\label{subsec:pbd-etapa}

Para la topología de la figura~\ref{crkt:sk-pasa-banda}, con $R_1 = R_3 = R$, $R_2 = 2R$ y $C_1 = C_2 = C$,
se tiene

\begin{align*}
  \omega_0 &= \frac{1}{R C} \\
  Q &= \frac{1}{3 - K} \\
  A_m &= \frac{K}{3 - K}
\end{align*}

con $K = 1 + R_b / R_a$.

\begin{figure}[!ht]
  \centering
  \begin{tikzpicture}[american]
  \node[op amp, noinv input up] (oa) at (6,0) {};

  \draw (0,0.5) node[ocirc, label=left:$V_i$] {} to[R, l=$R_1$] (1.8,0.5) node[circ] (n1) {}
    to[C, l=$C_2$] (3.6,0.5) node[circ] (n2) {} -- (oa.+);
  \draw (n1) to[R, l=$R_2$] ++(0,-2) node[ground] {};
  \draw (n2) to[R, l_=$R_3$] ++(0,-2) node[ground] {};
  \draw (oa.out) -- ++(0.6,0) node[circ] (out) {} to[short, -o] ++(1,0) node[right] {$V_o$};
  \draw (n1) to[C, l=$C_1$] ++(0,2) -| (out);

  % Red de ganancia K = 1 + Rb/Ra
  \draw (oa.-) -- ++(-0.5,0) coordinate (m) -- ++(0,-1.1) node[circ] (m2) {}
    to[R, l=$R_a$] ++(0,-1.8) node[ground] {};
  \draw (m2) to[R, l_=$R_b$] (out |- m2) -- (out);
\end{tikzpicture}

  \caption{Etapa Sallen-Key pasa banda de segundo orden con ganancia $K$.}
  \label{crkt:sk-pasa-banda}
\end{figure}

% TODO: desarrollar el cálculo de R, C, R_a y R_b. Si la ganancia A_m resultante no coincide con los 9 V/V pedidos,
% proponer un atenuador o una etapa de ganancia adicional y justificarlo.

\subsection{Valores comerciales}
\label{subsec:pbd-comerciales}

\begin{table}[!ht]
  \centering
  \begin{tabular}{|l|c|c|}
    \hline
    Componente & Calculado & Comercial \\
    \hline
    $R_1$ & & \\
    $R_2$ & & \\
    $R_3$ & & \\
    $C_1$ & & \\
    $C_2$ & & \\
    $R_a$ & & \\
    $R_b$ & & \\
    \hline
  \end{tabular}
  \caption{Valores calculados y comerciales del filtro pasa banda.}
  \label{tab:pbd-componentes}
\end{table}

\section{Simulación}
\label{sec:pbd-simulacion}

\subsection{Función de transferencia con elementos ideales}
\label{subsec:pbd-sim-ideal}

% TODO: diagrama de Bode (módulo y fase) con los valores calculados.

\subsection{Función de transferencia con valores comerciales}
\label{subsec:pbd-sim-real}

% TODO: repetir con los valores comerciales y con el modelo real del operacional; comparar f_o, BW y ganancia.

\section{Implementación y mediciones}
\label{sec:pbd-mediciones}

% TODO: si este filtro es uno de los dos implementados: armado, fotos en pictures/, tabla de barrido en frecuencia,
% captura del perfil de frecuencia (FFT) del osciloscopio y comparación medido vs. simulado. Si no se implementa,
% indicar que se realiza en el TP de filtros digitales.

\begin{table}[!ht]
  \centering
  \begin{tabular}{|c|c|c|c|c|}
    \hline
    $f$ [kHz] & $V_{in,pp}$ [V] & $V_{out,pp}$ [V] & $|H|$ medido [dB] & $|H|$ simulado [dB] \\
    \hline
     & & & & \\
     & & & & \\
     & & & & \\
    \hline
  \end{tabular}
  \caption{Relevamiento de la respuesta en frecuencia del filtro pasa banda.}
  \label{tab:pbd-barrido}
\end{table}
